\section{Conclusion and Future Work}\label{sec:conclusion-outlook}
This chapter answers the two research questions of Sec.~\ref{sec:research-question} and states what the study establishes. Sec.~\ref{sec:conclusion} gives the answers, and Sec.~\ref{sec:future-work} outlines the directions that follow from the results.

\subsection{Conclusion}\label{sec:conclusion}
The thesis examined whether GANs can generate TTF futures returns with the stylized facts of the observed data, namely heavy tails, volatility clustering, and asymmetry. In addition, it assessed whether price paths generated by recursive rollout reproduce the market prices of TTF options.

ForGAN and Fin-GAN were applied to daily TTF front-month futures returns. The generated returns were compared with the test split across training lengths, network capacities, condition window lengths, and input scaling methods. Each configuration was trained with five random seeds. Configuration 8 with the ForGAN baseline cost function was selected on the seed spread of the generated returns. The selected model generated price paths by recursive rollout. From these price paths, call options were priced on \num{20} valuation dates. The prices were compared with the market through the implied volatility. GBM and two bootstrap modes served as benchmarks.

The selected model reproduces the stylized facts in part. It learns heavy tails, with an excess kurtosis of \num{10.12} against \num{13.27} for the test split under configuration 8. The ForGAN baseline reproduces the sign of the asymmetry. The generated skewness of the ForGAN baseline is positive at \num{0.54}, below the \num{0.92} of the test split, while the $\mathrm{SR}^{*}$ term of Fin-GAN turns it negative. Configuration 8 comes closest to the volatility clustering of the test split at a lag of one day, but its autocorrelation of squared returns decays much more slowly at longer lags. No mode collapse occurred in any run.

The generated price paths do not yet reproduce the market prices of TTF options. The selected model has the highest implied volatility error at \num{0.256}, against \numrange{0.119}{0.149} for the benchmarks. The signed error is negative in every moneyness range. As a result, the deviation of the generated option prices from the market is a systematic underpricing. At the same time, the selected model reacts to the recent returns. The spread of its simulated returns correlates with the spread of the condition window at \num{0.758}, but stays below the market implied volatility on \num{17} of the \num{20} valuation dates.

For the setup of this thesis, \ref{rq:stylized} is answered in part and \ref{rq:pricing} negatively. The generated distribution is narrower than the observed returns. In addition, the recursive rollout moves $G$ outside the regime it was trained on.

\clearpage
\subsection{Future Work}\label{sec:future-work}

Although the generated prices do not match the market, the results show that the approach has the potential to capture the stylized facts of TTF futures returns and could be improved to better reproduce market option prices. The results point to several directions for further work.

\begin{enumerate}
    \item A longer training sample would test whether the tail weight of the generated returns increases with the amount of data. 
    \item The model selection and the pricing evaluation were both based on the test split. Repeating the selection on the validation split would show whether it produces the same configuration and cost function branch. It would also leave the test split as an independent basis for the pricing evaluation.
    \item GARCH-$t$ would add a benchmark that conditions on the preceding returns and would assess the results with a parametric model. 
    \item The seed spread of the one-step returns did not predict the seed spread of the option prices. A detailed analysis of the recursive rollout could clarify how errors accumulate and influence option pricing.
    \item The simulated spread of the selected model responds to the market state but stays below its level. Correcting the level of the generated volatility is left for future work. 
\end{enumerate}
These directions address the level of the generated volatility. This level is the gap between the response of the model to the market state and the prices quoted by the market.